\section{Objetivos}\label{cap:objetivos}

El objetivo principal de este Trabajo de Fin de Grado es el diseño y desarrollo de una aplicación móvil orientada a la Semana Santa, capaz de centralizar en una única plataforma la información de las distintas celebraciones de Semana Santa a nivel nacional, con independencia de la ciudad, cofradía o terminología empleada en cada una de ellas. En concreto, se busca ofrecer a los usuarios una herramienta que les permita consultar la información relevante de la Semana Santa de una ciudad y visualizar en tiempo real, mediante geolocalización, la ubicación, el recorrido y el estado de sus procesiones.

Para alcanzar este objetivo, se plantean los siguientes objetivos específicos:

\begin{itemize}
\item Diseñar dicha aplicación permitiendo a los usuarios elegir entre las distintas Semanas Santas dentro del territorio nacional.
\item Proporcionar acceso a la información general relacionada con la Semana Santa, incluyendo historia, cofradías o hermandades y elementos representativos.
\item Mostrar información detallada sobre los eventos, como nombre, descripción, fechas, horas e itinerarios.
\item Añadir funcionalidades de personalización, como la selección de eventos o elementos de interés y la recepción de notificaciones ante cambios o incidencias.
\item Implementar un sistema de visualización sobre mapa que permita representar las procesiones y su evolución en tiempo real.
\item Facilitar la navegación del usuario mediante la posibilidad de calcular rutas hacia ubicaciones concretas relacionadas con los eventos.
\item Desarrollar una solución accesible desde dispositivos móviles que garantice una experiencia de uso adecuada.
\item Diseñar una arquitectura escalable que permita la incorporación de nuevas ciudades y eventos sin necesidad de rediseñar el sistema.
\end{itemize}
